% !TEX program = context
% !TEX root = somerville-bubble-integrals.tex
% Run here: context somerville-bubble-integrals.tex
\input type-imp-somerville.mkxl
\input bubble-integrals.mkxl

% Somerville Regular: user coordinates, with the integral-slant axis at 100.
% Its source MGHT=700 maps to user MGHT=750 in type-imp-somerville.mkxl.
\defineSomervilleAxisTypeface{somerville-bubble-test}
  {wdth=100,wght=400,MGHT=750,opsz=12,STYA=5,STYB=5,ARLN=100,ARHD=100,INSL=100}
  {wdth=100,wght=400,opsz=12}
% Override the general typescript's goodie list for this test: somerville.lfg only.
\definefontfeature[somerville-bubble-test-math-axes]
  [axis={wdth=100,wght=400,MGHT=750,opsz=12,STYA=5,STYB=5,ARLN=100,ARHD=100,INSL=100},
   goodies={somerville}]
\setupbodyfont[somerville-bubble-test,12pt]
\setuppapersize[A4]
\setuplayout[backspace=18mm,topspace=16mm,width=middle,height=middle,
  header=0pt,footer=9mm]
\setupinteraction[state=start,title={Somerville: contour-aware integral and parenthesis spacing}]
\setuphead[chapter][style=\bfb,after={\blank[medium]}]
\setuphead[section][style=\bf,after={\blank[small]}]
\newbox\BubbleReferenceBox
\newbox\BubbleZeroBox

% Native ConTeXt operator, with an exact MATH variant for reproducible tests.
% Arguments: size index, limits/nolimits, upper limit, lower limit.
\unexpanded\def\TestFullIntegral#1#2#3#4#5{%
  \Uoperator #2 class \mathintegralcode
  \Umathchar \mathintegralcode \fam
    \getmathvariant\numexpr#1+1\relax "222B\relax
    {\doifnotempty{#3}{\mathstrut#3}}%
    {\doifnotempty{#4}{\mathstrut#4}}%
    {#5}\Uright rightclass \mathintegralcode .%
}
\unexpanded\def\TestIntegral#1#2#3#4{\TestFullIntegral{#1}{#2}{#3}{#4}{}}

\unexpanded\def\IntegralTest#1#2#3{%
  \startsection[title={#1}]
  \noindent Ordinary spacing:
  \startformula
    #2#3
  \stopformula
  \noindent Profile spacing:
  \startformula
    \bubbleintegral[trace=yes]{#2}{#3}
  \stopformula
  \stopsection
}

\starttext
\ctxlua{userdata.bubbleintegrals.selftest()}
\startchapter[title={Integral spacing by shape}]
Somerville Regular, 12pt optical design, integral slant 100.
The second version in each pair uses the actual outlines of the integral,
its limits, and the integrand. The first version uses normal ConTeXt spacing.

\IntegralTest{Short integrand}{\int[top=1,bottom=0]{}}{f(x)\,dx}
\IntegralTest{Fraction with a wide denominator}
  {\TestIntegral{8}{nolimits}{}{}}
  {\frac{1}{1+x+x^2+x^3}\,dx}
\IntegralTest{Fraction with a wide numerator}
  {\TestIntegral{8}{nolimits}{}{}}
  {\frac{1+x+x^2+x^3}{1}\,dx}
\stopchapter

\startchapter[title={Limits and nested constructions}]
\IntegralTest{Side limits participate in the profile}
  {\TestIntegral{12}{nolimits}{c+d}{a+b}}
  {\frac{\sqrt{1+x^2}}{1+x}\,dx}
\IntegralTest{Stacked limits participate in the profile}
  {\TestIntegral{12}{limits}{c+d}{a+b}}
  {\left(\frac{1}{1+\frac{x}{1+x}}\right)^2\,dx}
\IntegralTest{Tall integral with side limits}
  {\TestIntegral{16}{nolimits}{1}{0}}
  {\frac{x^2}{1+x^4}\,dx}
\stopchapter

\startchapter[title={Inspecting the envelopes}]
Grey lines show the actual glyph outlines and rules. Red shows the integral's
right-facing padded profile; blue shows the integrand's left-facing padded
profile. The profiles are compared over their full shared height.

\startformula
  \bubbleint[size=12,trace=yes]{\frac{1}{1+x+x^2+x^3}\,dx}
\stopformula
\midaligned{\scale[sx=2,sy=2]{\showbubbleintegralprofiles}}
\blank[big]
\startformula
  \bubbleint[size=12,trace=yes]{\frac{1+x+x^2+x^3}{1}\,dx}
\stopformula
\midaligned{\scale[sx=2,sy=2]{\showbubbleintegralprofiles}}
\stopchapter

\startchapter[title={Using the module}]
\starttyping
\input bubble-integrals.mkxl
\bubbleint[lower=0,upper=1]{f(x)\,dx}
\bubbleint[size=12]{\frac{1}{1+x^2}\,dx}
\bubbleintegral{\int[method=vertical,top=1,bottom=0]{}}{f(x)\,dx}
\stoptyping
The integrand is explicitly grouped and the resulting construction is
unbreakable. Ordinary \type{\int} remains available. This prototype targets
left-to-right mathematics made from glyphs and rules.

\startsection[title={Style and nesting checks}]
Inline: \m{\bubbleint[lower=0,upper=1]{f(x)\,dx}=1}.
\blank
Script: \m{X^{\bubbleint[lower=0,upper=1]{f(x)\,dx}}}.
\blank
Inline fraction: \m{\bubbleint{\frac{1}{1+x^2}\,dx}}.
\blank
Script fraction: \m{X^{\bubbleint{\frac{1}{1+x^2}\,dx}}}.
\blank
Nested:
\startformula
\bubbleint[lower=0,upper=1]
  {\bubbleint[lower=0,upper=x]{f(x,y)\,dy}\,dx}
\stopformula
\stopsection
\startsection[title={Bounding-box comparison}]
The first fit below uses the complete rectangular ink bounds; the second uses
the profiles. Both request the same clearance and padding.
\startformula
\bubbleint[size=12,method=bbox,trace=yes]{\frac{1}{1+x+x^2+x^3}\,dx}
\stopformula
\startformula
\bubbleint[size=12,trace=yes]{\frac{1}{1+x+x^2+x^3}\,dx}
\stopformula
\ctxlua{
  local m = userdata.bubbleintegrals
  local previous = m.history[rawlen(m.history)-1]
  assert(m.last.origin < previous.origin, "Profile fit should improve on the bounding-box fit")
  for fontid in pairs(m.last.left.fonts) do
    local font = fonts.hashes.identifiers[fontid]
    local instance = tostring(font.properties.instance)
    logs.report("bubble integral test", "operator instance: " .. instance)
    assert(instance:find("insl100",1,true) or instance:find("insl=100",1,true),
      "Test integral must use the INSL=100 instance")
  end
}
\stopsection
\stopchapter

\startchapter[title={Adjusting the strength}]
\type{strength=0} uses ordinary ConTeXt spacing. A value between 0 and 1
interpolates the horizontal position; \type{strength=1} applies the full
profile fit. The integral and its limits are unchanged by this parameter.

\startsection[title={Native integral, without the module}]
\startformula
\setbox\BubbleReferenceBox\mathstylehbox
  {\indisplaymath\conditionaltrue
   \TestFullIntegral{12}{limits}{1}{0}{\frac{1}{1+x^2}\,dx}}%
\setbox\BubbleZeroBox\mathstylehbox
  {\indisplaymath\conditionaltrue
   \bubbleint[strength=0,limits=yes,lower=0,upper=1,size=12]{\frac{1}{1+x^2}\,dx}}%
\ctxlua{
  local a,b = tex.getbox(\number\BubbleReferenceBox),tex.getbox(\number\BubbleZeroBox)
  assert(a.width==b.width and a.height==b.height and a.depth==b.depth,
    "Strength zero must preserve native layout dimensions")
  logs.report("bubble integral test","strength zero matches native layout dimensions")
}
\TestFullIntegral{12}{limits}{1}{0}{\frac{1}{1+x^2}\,dx}
\stopformula
\stopsection
\startsection[title={Strength 0: ordinary spacing}]
\startformula
\bubbleint[strength=0,limits=yes,lower=0,upper=1,size=12]{\frac{1}{1+x^2}\,dx}
\stopformula
\stopsection
\startsection[title={Strength 0.5: halfway}]
\startformula
\bubbleint[strength=.5,limits=yes,lower=0,upper=1,size=12,trace=yes]{\frac{1}{1+x^2}\,dx}
\stopformula
\ctxlua{
  local t = userdata.bubbleintegrals.last
  assert(math.abs(t.origin-(t.ordinaryorigin+t.targetorigin)/2) <= 1,
    "Half-strength must interpolate the horizontal position")
}
\stopsection
\startsection[title={Strength 1: full profile spacing}]
\startformula
\bubbleint[strength=1,limits=yes,lower=0,upper=1,size=12,trace=yes]{\frac{1}{1+x^2}\,dx}
\stopformula
\stopsection
\stopchapter
% Append before \stoptext in somerville-bubble-integrals.tex.
% The module and Somerville test face are already loaded by that document.
\unexpanded\def\ParenthesisPairTest#1#2#3{%
  \startsection[title={#1}]
  \noindent Ordinary spacing:
  \startformula
    #2
  \stopformula
  \noindent Profile spacing:
  \startformula
    #3
  \stopformula
  \noindent\showbubbleparenthesismetrics\par
  \stopsection
}
\unexpanded\def\ParenthesisTest#1#2{%
  \ParenthesisPairTest{#1}
    {\left(#2\right)}
    {\bubbleparentheses[trace=yes]{#2}}%
}

\startchapter[title={Parentheses: short expressions}]
The upper version of every pair uses ordinary automatically sized
parentheses. The lower version fits both parentheses to the expression's
profiles. The font, expression, and mathematical style remain the same.
The small figures give the changes from native spacing at the two joins.
A negative change brings the pieces closer; a positive change opens the gap.

\ParenthesisTest{A single letter}{x}
\ParenthesisTest{Italic stems and a descender}{f(x)+g(y)}
\ParenthesisTest{Subscripts and superscripts}{y_0^2+x_i^{n+1}}
\stopchapter

\startchapter[title={Parentheses: asymmetric fractions}]
These examples put the widest part of the expression at different heights.
The left and right sides are fitted independently.

\ParenthesisTest{Wide denominator}{\frac{1}{1+x+x^2+x^3}}
\ParenthesisTest{Wide numerator}{\frac{1+x+x^2+x^3}{1}}
\ParenthesisTest{Different shapes on the two sides}
  {\frac{x}{1+x^2}+\frac{1+x+x^2}{x}}
\stopchapter

\startchapter[title={Parentheses: radicals and scripts}]
Radicals and scripts contribute their occupied shapes. Rectangular bounds
alone do not describe how they fit against a parenthesis.

\ParenthesisTest{A square root}{\sqrt{1+x^2}}
\ParenthesisTest{Radicals of different heights}
  {\sqrt{\frac{1+x}{1-x}}+\sqrt{x}}
\ParenthesisTest{Scripts on both sides of a letter}
  {{}^{a+b}_{i+j}f(x)^{c+d}_{k+l}}
\stopchapter

\startchapter[title={Parentheses: nested constructions}]
The third pair has fitted inner parentheses in both versions, so the
comparison isolates the fit of the outer pair.

\ParenthesisTest{Nested fractions}{\frac{1}{1+\frac{x}{1+x}}}
\ParenthesisTest{An inner ordinary pair}
  {\left(\frac{1+x}{1-x}\right)^2+\frac{1}{x}}
\ParenthesisTest{An inner fitted pair}
  {\bubbleparentheses{\frac{1+x}{1-x}}^2+\frac{1}{x}}
\stopchapter

\startchapter[title={Parentheses: matrices and operators}]
These expressions have several separated regions of ink. Empty regions
between rows and between limits should not become solid rectangles.

\ParenthesisTest{A two by two matrix}
  {\startmatrix
     \NC a \NC b^2 \NR
     \NC c_i \NC d \NR
   \stopmatrix}
\ParenthesisTest{A column with an unusually wide middle row}
  {\startmatrix
     \NC x \NR
     \NC 1+x+x^2+x^3 \NR
     \NC 1 \NR
   \stopmatrix}
\ParenthesisTest{An integral and its stacked limits}
  {\int[method=vertical,top=1,bottom=0,size=24pt]
     {\frac{1}{1+x^2}\,dx}}
\stopchapter

\startchapter[title={Parentheses: integrals and empty regions}]
The first comparison changes the parentheses inside a fitted integral.
The remaining cases distinguish invisible size-setting material from ink.

\ParenthesisPairTest{Fitted parentheses inside a fitted integral}
  {\bubbleint[size=12]{\left(\frac{1+x}{1-x}\right)^2\,dx}}
  {\bubbleint[size=12]{\bubbleparentheses{\frac{1+x}{1-x}}^2\,dx}}
\ParenthesisTest{A short letter with a tall invisible strut}
  {x\vphantom{\frac{\frac{1}{1+x}}{\frac{1}{1+x}}}}
\ParenthesisPairTest{Empty and phantom-only parentheses}
  {\left(\right)\qquad
   \left(\vphantom{\frac{1}{1+x^2}}\right)}
  {\bubbleparentheses{}\qquad
   \bubbleparentheses{\vphantom{\frac{1}{1+x^2}}}}
\stopchapter

\startchapter[title={Parentheses: inspecting both fits}]
The enlarged drawings show how both sides respond to the expression.
In each drawing, the grey outlines are the actual glyphs and rules;
the coloured profiles are the boundaries used to choose the spacing.

\startsection[title={The wide part is below the axis}]
\startformula
\bubbleparentheses[trace=yes]{\frac{1}{1+x+x^2+x^3}}
\stopformula
\midaligned{\scale[sx=2,sy=2]{\showbubbleparenthesisprofiles}}
\stopsection
\startsection[title={The wide part is above the axis}]
\startformula
\bubbleparentheses[trace=yes]{\frac{1+x+x^2+x^3}{1}}
\stopformula
\midaligned{\scale[sx=2,sy=2]{\showbubbleparenthesisprofiles}}
\stopsection
\stopchapter

\startchapter[title={Parentheses: adjusting the strength}]
Every row uses the same expression. Zero gives ordinary spacing and one
gives the complete profile fit. Intermediate values interpolate both sides.

\startsection[title={Native parentheses, without the module}]
\startformula
\left(\frac{1}{1+x+x^2+x^3}\right)
\stopformula
\stopsection
\startsection[title={Strength 0}]
\startformula
\bubbleparentheses[strength=0]{\frac{1}{1+x+x^2+x^3}}
\stopformula
\stopsection
\startsection[title={Strength 0.25}]
\startformula
\bubbleparentheses[strength=.25]{\frac{1}{1+x+x^2+x^3}}
\stopformula
\stopsection
\startsection[title={Strength 0.5}]
\startformula
\bubbleparentheses[strength=.5]{\frac{1}{1+x+x^2+x^3}}
\stopformula
\stopsection
\startsection[title={Strength 0.75}]
\startformula
\bubbleparentheses[strength=.75]{\frac{1}{1+x+x^2+x^3}}
\stopformula
\stopsection
\startsection[title={Strength 1}]
\startformula
\bubbleparentheses[strength=1]{\frac{1}{1+x+x^2+x^3}}
\stopformula
\ctxlua{
  local m = userdata.bubbleintegrals
  local full = m.lastparentheses
  for i=rawlen(m.parenthesishistory)-3,rawlen(m.parenthesishistory) do
    local t = m.parenthesishistory[i]
    assert(math.abs(t.bodyorigin-(t.ordinarybody+t.strength*(full.targetbody-t.ordinarybody))) <= 1
      and math.abs(t.rightorigin-(t.ordinaryright+t.strength*(full.targetright-t.ordinaryright))) <= 1,
      "Changing strength must interpolate both parenthesis positions")
  end
  logs.report("bubble parenthesis test","both sides interpolate correctly at .25, .5, .75, and 1")
}
\stopsection
\stopchapter

\unexpanded\def\ParenthesisContextTest#1#2#3{%
  \startsection[title={#1}]
  Native: \m{#2}.\crlf
  Fitted: \m{#3}.
  \stopsection
}
\startchapter[title={Parentheses: inline and script styles}]
These pairs check the surrounding spacing and the smaller styles used in
scripts and fractions. Both versions use the same mathematical style.

\setbox\BubbleReferenceBox\hbox{\m{\left(\frac{1}{1+x}\right)}}%
\setbox\BubbleZeroBox\hbox{\m{\bubbleparentheses[strength=0]{\frac{1}{1+x}}}}%
\ctxlua{
  local a,b = tex.getbox(\number\BubbleReferenceBox),tex.getbox(\number\BubbleZeroBox)
  assert(a.width==b.width and a.height==b.height and a.depth==b.depth,
    "Strength zero must preserve native parenthesis dimensions")
  logs.report("bubble parenthesis test","strength zero matches native width, height, and depth")
}
\ParenthesisContextTest{Function argument and a following sum}
  {f\left(x+y\right)+z}
  {f\bubbleparentheses{x+y}+z}
\ParenthesisContextTest{A named function}
  {\sin\left(x+y\right)}
  {\sin\bubbleparentheses{x+y}}
\ParenthesisContextTest{Parentheses inside a superscript}
  {X^{\left(\frac{1}{1+x}\right)}}
  {X^{\bubbleparentheses{\frac{1}{1+x}}}}
\ParenthesisContextTest{Parentheses inside a subscript}
  {X_{\left(i+j\right)}}
  {X_{\bubbleparentheses{i+j}}}
\ParenthesisContextTest{Parentheses in a numerator}
  {\frac{\left(x+y\right)}{1+x^2}}
  {\frac{\bubbleparentheses{x+y}}{1+x^2}}
\ParenthesisContextTest{A superscript on the complete pair}
  {\left(\frac{1}{1+x}\right)^2}
  {\bubbleparentheses{\frac{1}{1+x}}^2}
\ParenthesisContextTest{Two adjacent pairs}
  {\left(x\right)\left(y\right)}
  {\bubbleparentheses{x}\bubbleparentheses{y}}
\stopchapter

\startchapter[title={Parentheses: very tall delimiters}]
The invisible rule sets a height of 110pt and a depth of 90pt. Only the
letter is ink inside the parentheses. Native delimiter sizing, including
extensible pieces when needed, is preserved; fitting changes the horizontal
positions only. The two parentheses must also clear each other.

\ParenthesisTest{An expression with 200pt of vertical extent}
  {x\vphantom{\vrule width0pt height110pt depth90pt}}

\startsection[title={Usage}]
\starttyping
\bubbleparentheses{\frac{1}{1+x^2}}
\bubbleparentheses[strength=.5]{x+y}
\bubbleparentheses[strength=0]{x+y}
\bubbleparentheses[clearance=.08em,padding=.02em]{f(x)}
\showbubbleparenthesisprofiles
\stoptyping
The short name \type{\bubbleparen} is also available. As with the integral
macro, each fitted construction is unbreakable.
\stopsection
\stopchapter

\stoptext
